\chapter{Introduction to Econometrics}
\label{ch:intro}

\section{What is econometrics?}

The term \emph{econometrics} was introduced by Ragnar Frisch, one of the
founders of the Econometric Society, the first editor of \emph{Econometrica},
and a co-recipient of the first Nobel Memorial Prize in Economic Sciences in
1969.  Frisch emphasized that econometrics is not economic statistics,
economic theory, or mathematics in isolation.  Its strength comes from
combining all three.  In modern language, we might summarize the field as
\[
  \text{economic models}+\text{statistical methods}+\text{economic data}.
\]
This chapter draws in part on Chapter~1 of Hansen's \emph{Econometrics}
(2022).

Econometrics therefore deals with data, but so do statistics and machine
learning.  What is distinctive about econometrics?  We will answer that
question partly here and, more importantly, throughout the course.

\section{An econometric description of data}

Suppose we want to know whether obtaining a master's degree would increase
future salary.  One tempting comparison is:
\begin{itemize}
  \item Alice has a master's degree and earns \$12,000 per month;
  \item Bob entered the labor market after college and earns \$8,000 per
        month;
  \item therefore, a master's degree raises salary by \$4,000.
\end{itemize}
This conclusion is not reliable.  Alice and Bob may differ in age, gender,
major, college, ability, family background, or labor-market conditions.  Two
people are also far too small a sample from which to draw a general
conclusion.  Two natural responses are to collect more characteristics for
each person (more \emph{variables}) and to collect information from more
people (a larger \emph{sample}).

\subsection{Variables and observations}

We usually use uppercase Latin letters for variables.  Let $Y_i$ be person
$i$'s monthly salary, let $X_{1i}$ indicate whether the person has a master's
degree, and let $X_{2i}$ be age.  We collect the explanatory variables in a
column vector,
\[
  X_i=(X_{1i},X_{2i})'.
\]
Unless stated otherwise, vectors in these notes are column vectors; a prime
denotes transpose.

Suppose a survey contains $n$ participants.  The sample can be written as
\[
  \{(Y_i,X_i):i=1,\ldots,n\}.
\]
It has $n$ observations and three variables.  In a data matrix, observations
are rows and variables are columns.  The following small example makes this
concrete.

\begin{rcode}
set.seed(5280)                 # make the demonstration reproducible
n <- 4
Y  <- exp(rnorm(n))            # simulated monthly income
X1 <- rbinom(n, 1, 0.5)        # 1 = master's degree
X2 <- sample(22:30, n, replace = TRUE)

data_matrix <- cbind(Y, X1, X2)
data_frame  <- data.frame(income = Y, MA = X1, age = X2)
data_matrix
data_frame
\end{rcode}

Mathematically, the data matrix is just a matrix.  We will use vectors and
matrices both to represent data and to derive estimators.  The notation takes
some getting used to, but it turns many long calculations into short ones.

\subsection{A probabilistic perspective}

Econometrics treats the sample as random.  Before we ask person $i$ about
salary, the answer is unknown to us and can be represented as a draw from a
distribution.  After the survey, we observe a \emph{realization} of that
random variable.  For four observations, the random data matrix is
\begin{equation}
\label{eq:intro-random-matrix}
  \begin{pmatrix}
    Y_1&X_{11}&X_{21}\\
    Y_2&X_{12}&X_{22}\\
    Y_3&X_{13}&X_{23}\\
    Y_4&X_{14}&X_{24}
  \end{pmatrix}.
\end{equation}
Before sampling, its entries are random; afterward, the realized data are
fixed numbers.

A useful way to organize the logic is the following.
\begin{enumerate}
  \item Nature determines a population distribution---the ``dice of God.''
        Researchers do not know it.
  \item We specify how observations will be sampled.
  \item Before seeing the outcomes, we specify an \emph{estimator} or a
        \emph{test statistic}: a rule that maps any possible sample into an
        answer.  Because the rule is a function of random observations, the
        estimator is itself random before the data are observed.
  \item We observe one realization of the sample.
  \item We apply the rule to that realization and obtain an \emph{estimate}
        or a test decision.
\end{enumerate}
The distinction between an estimator and an estimate is important.  The
estimator is the random rule; the estimate is the number produced by the
observed data.  In casual discussion, people also use \emph{sample} and
\emph{data set} interchangeably, although the former can emphasize the random
object and the latter its realization.

\subsection{Types of data}

Data can be classified along several dimensions.

\paragraph{Observational and experimental data.}
Observational data record outcomes generated without researcher-controlled
treatment assignment: surveys, stock returns, GDP, and prices are familiar
examples.  Experimental data arise when a researcher deliberately randomizes
an intervention.  Randomization often makes causal interpretation much
easier, but many economically important interventions cannot be randomized
for practical, legal, or ethical reasons.  We study methods for both kinds of
data.

\paragraph{Cross-sectional, time-series, and panel data.}
Cross-sectional data contain many entities observed at one time, such as a
one-time survey of 1,000 graduates.  Time-series data follow one entity over
many periods, such as quarterly GDP.  Panel data follow many entities over
many periods, such as an annual household survey.  Much of the course uses
cross-sectional notation, but panel data become central when we study
difference-in-differences and event studies in
Chapter~\ref{ch:did-event-study}.  In a panel, observations for the same unit
over time are generally dependent; later chapters account for that dependence
explicitly rather than imposing i.i.d. sampling at the unit--time level.

\section{Causal inference and prediction}

Causal inference and prediction are two major goals of data analysis.
Econometrics puts particular emphasis on causal questions:
\begin{itemize}
  \item Does air pollution reduce labor productivity?
  \item Does smoking during pregnancy harm infant health?
  \item Does graduate school raise earnings?
  \item Does a higher minimum wage reduce employment?
\end{itemize}
Prediction instead asks for an accurate forecast, for example tomorrow's
stock price, next quarter's GDP, or which product a customer is likely to
buy.  A model may predict well without describing what would happen under an
intervention.  Conversely, a credible causal design need not be the best
forecasting model.

Controlled experiments provide especially transparent causal evidence, but
they are not always feasible.  Econometricians have therefore developed a
large toolkit for recovering causal effects from observational data.  This
course studies designs based on randomized experiments as well as selection
on observables, instrumental variables, regression discontinuity,
difference-in-differences, event studies, and modern machine-learning tools.
